
This work demonstrates that two simple statistical features---normalization layer type counts and parameter-mass ratio---achieve 88.89\% architecture family classification accuracy with $100\times$ faster checkpoint extraction and full interpretability compared to graph neural network approaches requiring 50+ hours implementation and GPU-intensive graph construction.

The central insight is that architecture families impose structural constraints observable as checkpoint fingerprints. CNNs use BatchNorm for spatial normalization and allocate parameters to convolutional kernels. Transformers use LayerNorm for token-wise normalization and allocate parameters to linear projection matrices. These are architectural paradigms directly visible in checkpoint metadata through layer names and tensor shapes.

Validation across five experiments demonstrates 88.89\% accuracy on held-out TIMM models, 83.3\% accuracy on edge case architectures with 1.7\% degradation, perfect scale invariance (CV = 0.00 across ResNet family spanning $5\times$ parameter range), and exceptionally strong inter-family separation (Cohen's d = 3.202, p < 0.001). Mechanistic validation confirms CNNs exclusively use BatchNorm (0\% violation) while Transformers predominantly use LayerNorm (14.29\% violation). Checkpoint-only extraction completes in 61 seconds with 0 MB GPU usage.

Principal limitation is complete failure on NormFree networks (0\% accuracy on NFNet) that replace normalization layers with scaled weight standardization. With absent normalization, classification relies solely on parameter-mass ratio, which is insufficient for these paradigms. This limitation affects rare architectures and has clear mechanistic cause. Small validation set (18 models) limits statistical power, though primary metrics exceed thresholds with comfortable margins.

The contribution challenges complexity paradigm in weight-space learning. When targeting specific classification tasks aligned with hand-crafted features guided by mechanistic understanding, simple statistical approaches match complex neural methods with better interpretability and efficiency. Architecture families are structural categories observable in checkpoint fingerprints, not merely behavioral categories requiring execution. This establishes foundation for interpretable weight-space analysis: extracting information through features that capture known architectural constraints rather than learning opaque representations through neural networks.
